% ============================================================
%  DobleA LAB — Módulo 1
%  De la inquietud del cliente al diseño de estudio
%  Miércoles 6 de mayo · 6–9 PM
% ============================================================
\documentclass[aspectratio=169, 12pt]{beamer}

% ── Encoding & language ────────────────────────────────────
\usepackage[utf8]{inputenc}
\usepackage[T1]{fontenc}
\usepackage[spanish]{babel}
\usepackage{microtype}

% ── Fonts ──────────────────────────────────────────────────
\usepackage{lmodern}

% ── Colors (DobleA palette) ────────────────────────────────
\usepackage{xcolor}
\definecolor{daPurpleDeep}{HTML}{3B0F8C}
\definecolor{daPurpleMid}{HTML}{6B21C8}
\definecolor{daPurpleLight}{HTML}{8B5CF6}
\definecolor{daMagenta}{HTML}{E040A0}
\definecolor{daDark}{HTML}{0A0A0A}
\definecolor{daLightBg}{HTML}{F0EDE8}
\definecolor{daTextMid}{HTML}{444444}

% ── Beamer theme base ──────────────────────────────────────
\usetheme{default}
\usecolortheme{default}
\setbeamertemplate{navigation symbols}{}
\setbeamertemplate{itemize item}{\color{daMagenta}\textbullet}
\setbeamertemplate{itemize subitem}{\color{daPurpleLight}$\circ$}

% ── Background ─────────────────────────────────────────────
\setbeamercolor{background canvas}{bg=daDark}
\setbeamercolor{normal text}{fg=white}
\setbeamercolor{frametitle}{fg=white, bg=daPurpleDeep}
\setbeamercolor{title}{fg=white}
\setbeamercolor{subtitle}{fg=daPurpleLight}
\setbeamercolor{author}{fg=daLightBg}
\setbeamercolor{date}{fg=daMagenta}
\setbeamercolor{institute}{fg=daLightBg}
\setbeamercolor{block title}{bg=daPurpleMid, fg=white}
\setbeamercolor{block body}{bg=daPurpleDeep!40!black, fg=white}
\setbeamercolor{alerted text}{fg=daMagenta}

% ── Frametitle style ───────────────────────────────────────
\setbeamertemplate{frametitle}{%
  \nointerlineskip
  \begin{beamercolorbox}[wd=\paperwidth, ht=0.55cm, dp=0.25cm, leftskip=0.8cm]{frametitle}%
    \usebeamerfont{frametitle}\insertframetitle%
    \ifx\insertframesubtitle\empty\else
      \hspace{0.5em}{\color{daMagenta}\small\textit{— \insertframesubtitle}}%
    \fi
  \end{beamercolorbox}%
  \vspace{-0.15cm}%
  {\color{daMagenta}\rule{\paperwidth}{1.2pt}}%
}

% ── Footer ─────────────────────────────────────────────────
\setbeamertemplate{footline}{%
  {\color{daPurpleDeep!60!black}\rule{\paperwidth}{0.5pt}}%
  \begin{beamercolorbox}[wd=\paperwidth, ht=0.45cm, dp=0.15cm]{normal text}%
    \hspace{0.6cm}%
    {\color{daMagenta}\tiny DobleA LAB}%
    \hfill%
    {\color{daLightBg!60!black}\tiny Módulo 1 · De la inquietud del cliente al diseño de estudio}%
    \hfill%
    {\color{daPurpleLight}\tiny\insertframenumber/\inserttotalframenumber\hspace{0.5cm}}%
  \end{beamercolorbox}%
}

% ── Title page template ────────────────────────────────────
\setbeamertemplate{title page}{%
  \begin{tikzpicture}[remember picture, overlay]
    % Gradient background strip
    \fill[daPurpleDeep] (current page.north west) rectangle ([yshift=-0.45\paperheight]current page.north east);
    \fill[daPurpleMid!60!daDark] ([yshift=-0.45\paperheight]current page.north west) rectangle (current page.south east);
    % Magenta accent line
    \draw[daMagenta, line width=2pt] ([yshift=-0.45\paperheight]current page.north west) -- ([yshift=-0.45\paperheight]current page.north east);
  \end{tikzpicture}%
  \vspace{0.8cm}
  \begin{center}
    {\color{daMagenta}\footnotesize\textbf{\MakeUppercase{DobleA LAB · Taller de Investigación Aplicada}}}\\[0.3cm]
    {\color{white}\LARGE\textbf{\inserttitle}}\\[0.3cm]
    {\color{daPurpleLight}\small\insertsubtitle}\\[0.6cm]
    {\color{daLightBg}\footnotesize\insertauthor}\\[0.15cm]
    {\color{daMagenta}\footnotesize\insertdate}
  \end{center}
}

% ── TikZ for title bg ──────────────────────────────────────
\usepackage{tikz}

% ── Metadata ───────────────────────────────────────────────
\title{De la inquietud del cliente\\al diseño de estudio}
\subtitle{Módulo 1}
\author{Isidora Aninat \and Lucía Argote \and Pablo Argote}
\date{Miércoles 6 de mayo · 6--9 PM}
\institute{DobleA LAB}

% ── Helper command: session header ─────────────────────────
\newcommand{\sesion}[2]{%
  \begin{beamercolorbox}[wd=\linewidth, ht=0.4cm, dp=0.15cm, leftskip=0.3cm, rounded=true]{block title}%
    {\small\textbf{Sesión #1} \hfill \textit{\footnotesize #2}}%
  \end{beamercolorbox}\vspace{0.15cm}%
}

% ============================================================
\begin{document}

% ── Title slide ────────────────────────────────────────────
\begin{frame}[plain]
  \titlepage
\end{frame}

% ── Overview ───────────────────────────────────────────────
\begin{frame}{Contenidos del módulo}
  \begin{columns}[T]
    \begin{column}{0.32\textwidth}
      \sesion{1}{50 min}
      {\small
      \begin{itemize}
        \item Del problema a la pregunta de investigación
        \item Objetivos y conceptos medibles
      \end{itemize}}
    \end{column}
    \begin{column}{0.32\textwidth}
      \sesion{2}{50 min}
      {\small
      \begin{itemize}
        \item Estrategias metodológicas
        \item Técnicas cualitativas
      \end{itemize}}
    \end{column}
    \begin{column}{0.32\textwidth}
      \sesion{3}{50 min}
      {\small
      \begin{itemize}
        \item Diseño del estudio
        \item Muestreo cualitativo
        \item Ejercicio práctico
      \end{itemize}}
    \end{column}
  \end{columns}
\end{frame}

% ============================================================
%  SESIÓN 1
% ============================================================
\begin{frame}[plain]
  \begin{tikzpicture}[remember picture, overlay]
    \fill[daPurpleMid] (current page.north west) rectangle (current page.south east);
    \node[anchor=center] at (current page.center) {%
      \color{white}\Huge\textbf{Sesión 1}%
    };
    \node[anchor=center, yshift=-1.2cm] at (current page.center) {%
      \color{daMagenta}\large Del problema al diseño%
    };
  \end{tikzpicture}
\end{frame}

\begin{frame}{Recoger la necesidad del cliente}
  \begin{block}{Punto de partida}
    Todo estudio comienza con una \alert{inquietud}: algo que no se sabe y que importa saber.
  \end{block}
  \vspace{0.4cm}
  \begin{itemize}
    \item Escuchar y estructurar la necesidad de información del cliente.
    \item Distinguir entre \textit{problema de negocio} y \textit{pregunta de investigación}.
    \item Identificar qué decisión se quiere informar con el estudio.
  \end{itemize}
\end{frame}

\begin{frame}{De la inquietud a la pregunta de investigación}
  \begin{columns}[T]
    \begin{column}{0.48\textwidth}
      \begin{block}{Inquietud del cliente}
        \small``No sabemos por qué nuestros clientes se van después del primer año.''
      \end{block}
    \end{column}
    \begin{column}{0.48\textwidth}
      \begin{block}{Pregunta de investigación}
        \small¿Qué factores están asociados a la no renovación de contrato entre clientes de primer año?
      \end{block}
    \end{column}
  \end{columns}
  \vspace{0.5cm}
  \begin{itemize}
    \item La pregunta debe ser \alert{específica}, \alert{investigable} y \alert{relevante}.
    \item Definir objetivos generales y específicos.
    \item Operacionalizar conceptos: ¿qué vamos a medir exactamente?
  \end{itemize}
\end{frame}

\begin{frame}{Preguntas que no funcionan}
  \framesubtitle{Anti-ejemplos}
  \small
  \begin{columns}[T]
    \begin{column}{0.48\textwidth}
      \textbf{\color{daMagenta}Mal formulada}
      \begin{itemize}
        \item ``¿Es bueno nuestro servicio?''\\
              {\color{daTextMid}\tiny Demasiado vaga; no se puede operacionalizar.}
        \item ``¿Por qué la gente no nos compra?''\\
              {\color{daTextMid}\tiny Asume que hay una sola razón; no define población.}
        \item ``¿Cómo mejorar las ventas?''\\
              {\color{daTextMid}\tiny Es un objetivo de negocio, no una pregunta investigable.}
        \item ``¿Los millennials prefieren nuestra marca?''\\
              {\color{daTextMid}\tiny Contiene sesgo de confirmación; categoría imprecisa.}
      \end{itemize}
    \end{column}
    \begin{column}{0.48\textwidth}
      \textbf{\color{daPurpleLight}Bien formulada}
      \begin{itemize}
        \item ``¿Cómo evalúan la experiencia de postventa los clientes que renovaron vs.\ los que no?''
        \item ``¿Qué factores inciden en la decisión de no compra entre mujeres 25--40 en la RM?''
        \item ``¿Qué barreras perciben los equipos de venta para cerrar negocios en el segmento PYME?''
        \item ``¿Qué atributos de marca asocian espontáneamente los consumidores 20--35 a nuestra categoría?''
      \end{itemize}
    \end{column}
  \end{columns}
\end{frame}

\begin{frame}{¿Por qué fallan?}
  \framesubtitle{Errores frecuentes al formular preguntas}
  \begin{enumerate}
    \item \textbf{Vaguedad:} no define qué se quiere saber ni sobre quién.
    \item \textbf{Sesgo implícito:} la pregunta ya contiene la respuesta que se espera.
    \item \textbf{Confundir objetivo de negocio con pregunta de investigación:}\\
          {\small ``aumentar ventas'' no es algo que se investiga, es algo que se \textit{decide} a partir de la investigación.}
    \item \textbf{No investigable:} no se puede responder con datos recolectables\\
          {\small (ej.\ ``¿Qué va a pasar con el mercado en 10 años?'').}
    \item \textbf{Demasiado amplia:} intenta responder todo a la vez.
  \end{enumerate}
  \vspace{0.3cm}
  \begin{alertblock}{Test rápido}
    Si no puedes imaginar \textit{qué dato} responde tu pregunta, la pregunta no está lista.
  \end{alertblock}
\end{frame}

\begin{frame}{Definir conceptos medibles}
  \begin{itemize}
    \item Un concepto abstracto (\textit{satisfacción}, \textit{confianza}) debe traducirse en indicadores observables.
    \item \textbf{Dimensiones} → \textbf{Indicadores} → \textbf{Ítems de medición}
  \end{itemize}
  \vspace{0.4cm}
  \begin{exampleblock}{Ejemplo: ``Satisfacción con el servicio''}
    \small
    \textbf{Dimensiones:} calidad percibida · tiempo de respuesta · trato recibido\\
    \textbf{Indicadores:} NPS · tiempos de resolución · escala de valoración\\
    \textbf{Ítems:} ``¿Cuán probable es que recomiende este servicio? (0–10)''
  \end{exampleblock}
\end{frame}

% ── Descanso ───────────────────────────────────────────────
\begin{frame}[plain]
  \begin{tikzpicture}[remember picture, overlay]
    \fill[daDark] (current page.north west) rectangle (current page.south east);
    \draw[daMagenta, line width=1pt] ([xshift=2cm, yshift=-2cm]current page.north west) rectangle ([xshift=-2cm, yshift=2cm]current page.south east);
    \node[anchor=center] at (current page.center) {%
      \color{daPurpleLight}\large\textit{— break 10 min —}%
    };
  \end{tikzpicture}
\end{frame}

% ============================================================
%  SESIÓN 2
% ============================================================
\begin{frame}[plain]
  \begin{tikzpicture}[remember picture, overlay]
    \fill[daPurpleMid] (current page.north west) rectangle (current page.south east);
    \node[anchor=center] at (current page.center) {%
      \color{white}\Huge\textbf{Sesión 2}%
    };
    \node[anchor=center, yshift=-1.2cm] at (current page.center) {%
      \color{daMagenta}\large Estrategias metodológicas%
    };
  \end{tikzpicture}
\end{frame}

\begin{frame}{¿Cuándo usar qué metodología?}
  \begin{columns}[T]
    \begin{column}{0.31\textwidth}
      \begin{block}{\centering Cualitativa}
        \small Comprender \textit{por qué} y \textit{cómo}.\\ Explorar significados.\\ Hipótesis emergentes.
      \end{block}
    \end{column}
    \begin{column}{0.31\textwidth}
      \begin{block}{\centering Cuantitativa}
        \small Medir \textit{cuánto} y \textit{con qué frecuencia}.\\ Generalizar a poblaciones.\\ Probar hipótesis.
      \end{block}
    \end{column}
    \begin{column}{0.31\textwidth}
      \begin{block}{\centering Mixta}
        \small Combinar profundidad y representatividad.\\ Triangulación de resultados.
      \end{block}
    \end{column}
  \end{columns}
  \vspace{0.4cm}
  {\small La elección depende de la \alert{pregunta de investigación}, no de la preferencia del investigador.}
\end{frame}

\begin{frame}{Técnicas cualitativas}
  \begin{itemize}
    \item \textbf{Entrevista en profundidad:} diálogo individual, relato detallado, alta flexibilidad.
    \begin{itemize}
      \item Ideal para trayectorias, experiencias personales, temas sensibles.
    \end{itemize}
    \item \textbf{Grupos focales:} interacción grupal, construcción colectiva de sentido.
    \begin{itemize}
      \item Ideal para normas sociales, discursos compartidos, evaluación de productos.
    \end{itemize}
    \item \textbf{Etnografía / observación:} presencia en el campo, comportamiento natural.
    \begin{itemize}
      \item Ideal para prácticas cotidianas, entornos específicos.
    \end{itemize}
  \end{itemize}
\end{frame}

\begin{frame}{Criterios de elección de técnica}
  \begin{block}{Preguntas para elegir}
    \begin{itemize}
      \item ¿Qué tipo de dato necesito? (narrativo · numérico · conductual)
      \item ¿Cuánta profundidad vs.\ amplitud requiero?
      \item ¿Qué recursos y tiempo tengo disponibles?
      \item ¿Qué tan sensible es el tema para los participantes?
    \end{itemize}
  \end{block}
\end{frame}

% ── Descanso ───────────────────────────────────────────────
\begin{frame}[plain]
  \begin{tikzpicture}[remember picture, overlay]
    \fill[daDark] (current page.north west) rectangle (current page.south east);
    \draw[daMagenta, line width=1pt] ([xshift=2cm, yshift=-2cm]current page.north west) rectangle ([xshift=-2cm, yshift=2cm]current page.south east);
    \node[anchor=center] at (current page.center) {%
      \color{daPurpleLight}\large\textit{— break 10 min —}%
    };
  \end{tikzpicture}
\end{frame}

% ============================================================
%  SESIÓN 3
% ============================================================
\begin{frame}[plain]
  \begin{tikzpicture}[remember picture, overlay]
    \fill[daPurpleMid] (current page.north west) rectangle (current page.south east);
    \node[anchor=center] at (current page.center) {%
      \color{white}\Huge\textbf{Sesión 3}%
    };
    \node[anchor=center, yshift=-1.2cm] at (current page.center) {%
      \color{daMagenta}\large Diseño del estudio y muestreo%
    };
  \end{tikzpicture}
\end{frame}

\begin{frame}{Definir la estrategia de estudio}
  \begin{itemize}
    \item \textbf{Tipo de estudio:} exploratorio · descriptivo · explicativo · evaluativo
    \item \textbf{Alcance:} transversal vs.\ longitudinal
    \item \textbf{Unidad de análisis:} persona · hogar · organización · evento
    \item \textbf{Población objetivo:} ¿a quiénes queremos conocer?
  \end{itemize}
  \vspace{0.3cm}
  \begin{alertblock}{Regla de oro}
    El diseño responde a la pregunta: no al revés.
  \end{alertblock}
\end{frame}

\begin{frame}{Muestreo en estudios cualitativos}
  \begin{itemize}
    \item El objetivo no es representatividad estadística sino \alert{representatividad tipológica}.
    \item \textbf{Criterios de selección:} relevancia · variedad · acceso
    \item \textbf{Saturación teórica:} cuando los casos nuevos ya no agregan información nueva.
    \item Tamaño típico: 8--20 entrevistas (depende de la homogeneidad del grupo).
  \end{itemize}
  \vspace{0.3cm}
  \begin{block}{Tipos de muestreo}
    Intencional · bola de nieve · máxima variación · casos típicos
  \end{block}
\end{frame}

\begin{frame}{Ejercicio práctico}
  \begin{block}{Instrucción}
    En grupos de 3, tomen el siguiente escenario y desarrollen:
  \end{block}
  \vspace{0.3cm}
  \textbf{Escenario:} Una ONG quiere saber por qué la participación en sus talleres comunitarios bajó un 40\% en el último año.
  \vspace{0.3cm}
  \begin{enumerate}
    \item Formula una \textbf{pregunta de investigación} precisa.
    \item Elige una \textbf{estrategia metodológica} y justifícala.
    \item Define la \textbf{técnica} y los \textbf{criterios de selección} de participantes.
    \item ¿Cuántos casos necesitarías y por qué?
  \end{enumerate}
\end{frame}

% ── Cierre ─────────────────────────────────────────────────
\begin{frame}{Cierre del módulo}
  \begin{columns}[T]
    \begin{column}{0.48\textwidth}
      \textbf{\color{daMagenta}Lo que aprendimos hoy}
      \begin{itemize}
        \small
        \item Traducir necesidades en preguntas investigables.
        \item Operacionalizar conceptos.
        \item Elegir estrategia y técnica metodológica.
        \item Diseñar un muestreo cualitativo.
      \end{itemize}
    \end{column}
    \begin{column}{0.48\textwidth}
      \textbf{\color{daPurpleLight}Próximo módulo}
      \begin{itemize}
        \small
        \item Diseño de instrumentos de medición.
        \item Calidad en la medición.
        \item Miércoles 13 de mayo · 6--9 PM
      \end{itemize}
    \end{column}
  \end{columns}
  \vspace{0.5cm}
  \begin{center}
    {\color{daMagenta}\rule{6cm}{0.8pt}}\\[0.3cm]
    {\small\color{daLightBg!70!black} Preguntas · Comentarios · Dudas}
  \end{center}
\end{frame}

\end{document}
